% !TeX root = ../../main.tex

\subsubsection{Summary of Neurons and Synapses for a \( 4 \times 4 \) Sudoku Grid}

For a \( 4 \times 4 \) Sudoku, 64 neurons are needed to represent the state of the board.

In addition to this, 64 stimulating input neurons/spikes are used.
Each of these is connected to one neuron, meaning 64 stimulating synapses.

There is 1 additional stimulating input spiking neuron connected to all neurons.
The synapse weights for this neuron are adjusted to 0 or 63 depending on if the corresponding clue was given or not.
This leads to a maximum of \( 4^2 = 16 \) \quoteapprox{active} synapses with weight \( ≠ 0 \), if each field has a clue.

Constraint Synapses:
\begin{itemize}
    \item The constraint for one number per field means for each of the \( 4^2 = 16 \) fields, inhibitory synapses are needed to connect each of the 4 neurons representing each number to the other 3.
    \subitem This means \( 4^2 \cdot 4 \cdot 3 = 192 \) synapses.

    \item The constraint for no repeat numbers in each row means for each of the \( [4 \ \text{numbers} \cdot 4 \ \text{rows} = 16] \) neurons in rows,
    inhibitory synapses are needed connecting each of the 4 neurons to the other 3.
    \subitem This means \( 4^2 \cdot 4 \cdot 3 = 192 \) synapses.

    \item The constraint for no repeat numbers in each column similarly means each of the \( [4 \ \text{numbers} \cdot 4 \ \text{rows} = 16] \) neurons columns,
    inhibitory synapses connect each of the 4 neurons to the other 3.
    \subitem This means \( 4^2 \cdot 4 \cdot 3 = 192 \) synapses.

    \item The last constraint calls for no repeat numbers in each of the \( [2^2 = 4] \) blocks with shape: \( (2 x 2) \), as the rules of Sudoku dictate.
    This constraint applies to each number, so there are in total \( [2^3 = 8] \) blocks of 4 neurons.
    Inhibitory synapses must connect each of the 4 neurons in each block (with same number) to the other 3.
    \subitem This means \( 2^3 \cdot 4 \cdot 3 = 96 \) synapses.
\end{itemize}

Concluding the summary:

There are 64 main neurons doing the calculations.
Counting the \( [64 + 1 = 65] \) stimulating input spike neurons leads to a total of:
\[129 \ \text{neurons.}\]
The calculations considering \( [64 + 16 = 80] \) stimulating neurons and 672 inhibiting neurons lead to a total of:
\[752 \ \text{synapses.}\]